\section{Ordered cache simulation on 100 GPT-OSS problems}
\label{sec:oss-cache-100}

The first completed 100-problem cache experiment tests fixed positive
identity bias $+1$ on GPT-OSS, with one baseline and one guided generation
per held-out problem. It uses the frozen OSS policy, matched prompts and
per-request seeds, temperature 0.6, top-$p$ 0.95, a 32,768-token output
limit, and synchronous single-sequence execution. The deterministic
round-robin selection contains 9 AIME24, 9 AIME25, 12 AMC23, 24 MATH-500,
23 Minerva and 23 OlympiadBench problems. Dataset sizes are unequal because
the smaller held-out pools are exhausted. This is a separate generation
experiment from the full-diagnostics run in
Section~\ref{sec:oss-cache-diagnostics}; their aggregate effects need not agree.

\paragraph{Quality and generation length.}
The original generation scorer gives 53/100 correct baseline answers and
48/100 guided answers. Applying the report's offline extracted-answer
equivalence contract to both conditions gives 76/100 and 75/100 instead.
The rescored change is $-1.00$ percentage point, with a paired,
dataset-stratified problem-bootstrap 95\% interval $[-7.00,+5.00]$
(5,000 resamples, seed 42). Six answers change from wrong to right and
seven from right to wrong. Neither accuracy preservation nor a systematic
accuracy loss is established. The large scoring discrepancy reinforces
why the original and rescored results must remain explicitly distinguished.

Mean generated tokens increase from 4,418.86 to 5,392.95 ($+22.04\%$), and
token-limit hits rise from 2 to 6. The largest length increase is on AIME24
(8,957.0 to 16,580.1 mean tokens). Rescored correct counts change from 4 to 5
on AIME24, 8 to 4 on AIME25, 12 to 11 on AMC23, 23 to 23 on MATH-500,
13 to 13 on Minerva, and 16 to 19 on OlympiadBench. These small
benchmark cohorts describe the composition of the aggregate effect rather
than establishing general model-level performance differences.

\paragraph{Cache measurement and controls.}
Ordered top-4 routing sets are recorded at all 24 layers, excluding all
prefill chunks and the first output token predicted by prefill. Every
trace is checked against its generation length and source hash. The
simulation starts with a cold cache for each answer and uses equal
capacity at each layer. It compares baseline LRU, baseline with
calibration-frequency pins, guided routing with target pins, and guided
LRU. Four experts are pinned at each of the four policy layers; other
layers use LRU. Initial pin loads count, including pins never requested.
Simultaneous top-$k$ requests cannot evict each other, and sorted expert
IDs break simultaneous recency ties. Frequency pins are chosen from
calibration statistics, not evaluation traces.

\begin{table}[ht]
\centering\small
\caption{GPT-OSS, 100 held-out problems: simulated decode-only expert loads. Every response starts cold; initial pin loads are included. Slots are per layer.}
\label{tab:oss-cache-100-loads}
\begin{tabular}{llrrr}
\toprule
Policy & Slots & Loads/answer & Loads/token & Hit rate (\%) \\
\midrule
Baseline LRU & 8 & 134,007.0 & 30.33 & 68.40 \\
Baseline LRU & 16 & 40,079.7 & 9.07 & 90.55 \\
Baseline LRU & 32 & 722.8 & 0.16 & 99.83 \\
Baseline frequency pins & 8 & 134,747.1 & 30.50 & 68.23 \\
Baseline frequency pins & 16 & 40,068.2 & 9.07 & 90.56 \\
Baseline frequency pins & 32 & 722.8 & 0.16 & 99.83 \\
Guided target pins & 8 & 162,649.3 & 30.17 & 68.58 \\
Guided target pins & 16 & 46,092.0 & 8.55 & 91.10 \\
Guided target pins & 32 & 717.1 & 0.13 & 99.86 \\
Guided LRU & 8 & 158,902.8 & 29.47 & 69.30 \\
Guided LRU & 16 & 45,581.3 & 8.45 & 91.19 \\
Guided LRU & 32 & 717.1 & 0.13 & 99.86 \\
\bottomrule
\end{tabular}
\end{table}


\paragraph{Findings.}
At eight slots per layer, guided target pinning increases mean loads per
answer from 134,007.03 to 162,649.30 ($+21.37\%$), while loads per decode
token change by only $-0.55\%$. At 16 slots, loads per token fall from
9.07 to 8.55 ($-5.77\%$), but loads per answer increase from 40,079.70 to
46,092.03 ($+15.00\%$). Thus improved per-token locality is insufficient
to offset the additional generated tokens in the observed sample.

Paired dataset-stratified problem resampling of the trace-derived totals
(5,000 resamples, Python random seed 42, percentile intervals) gives a
95\% interval of $[-9.06,-2.56]\%$ for the 16-slot per-token change and
$[-13.33,+53.22]\%$ for its per-answer change. At eight slots, the
corresponding intervals are $[-2.74,+1.66]\%$ and $[-8.92,+62.53]\%$.
Intervals are conditional on saved generations and the frozen policy,
with no adjustment for multiple comparisons. The per-token reduction at
16 slots is supported under this resampling procedure; a population-level
increase or decrease in total loads per answer is not resolved.

On the same guided traces, ordinary LRU uses fewer loads than target
pinning: 158,902.83 versus 162,649.30 per answer at eight slots, and
45,581.32 versus 46,092.03 at 16 slots. Pinning therefore adds overhead
relative to adaptive caching for these trajectories. The frequency-pinned
baseline is close to baseline LRU (134,747.10 and 40,068.22 loads per
answer at eight and 16 slots). At 32 slots all 32 OSS experts fit at each
layer: the small decrease from 722.75 to 717.10 loads per answer concerns
initial loading of distinct experts, not reduced eviction under memory
pressure. Its relative change is $-0.78\%$ ($[-1.05,-0.50]\%$).

\paragraph{Interpretation and limits.}
This run provides evidence of improved per-token locality at one
constrained cache budget, but no demonstrated total-loading or speed
advantage for the tested guided-and-pinned policy. A modest quality loss
could be acceptable if compensated by sufficient speed; this experiment
does not establish that compensation. Output length and loads per answer
must accompany hit rates when judging an efficiency trade-off.

These are simulated equal-sized expert loads, not measured SSD reads or
bytes. Capture derives top-$k$ identities from router-hook logits rather
than observing native dispatch, so tied scores can differ. The simulation
omits prefill warming, cross-answer reuse, batching, prefetch and transfer
costs. All expert weights remain resident during generation, and recorder
overhead prevents interpreting captured generation time as offloaded
inference latency. The result is limited to one fixed policy, one sampled
trajectory per condition per problem, and this 100-problem cohort.
